\bookchapter{Capitolo 33: Prospettive future — il prossimo viaggio del telefono}

\chapterintro{Il telefono non scomparirà dall'oggi al domani, ma la sua “posizione egemone” potrebbe essere erosa da una nuova generazione di dispositivi. In questo capitolo parliamo di tre cose: che aspetto avranno i dispositivi di prossima generazione, come il telefono si “invisibilizzerà” diventando la nuova pietra angolare dell'identità, e le questioni etiche che dobbiamo cominciare a pensare in anticipo.}

\booksection{AR/VR e calcolo spaziale: la “seconda scrivania” oltre il telefono}

La competizione sui dispositivi di prossima generazione è già iniziata, e quelli con maggiori probabilità di erodere il tempo passato sul telefono sono i dispositivi AR/VR e di calcolo spaziale:

\begin{itemize}
  \item \textbf{Il seguito di Vision Pro}: il lancio di Apple Vision Pro ha segnato la concreta materializzazione del concetto di "calcolo spaziale", ma prezzo (da 29999 yuan) e peso della prima generazione restano soglie d'ingresso. Si prevede che nel 2026-2027 uscirà una versione Vision più leggera ed economica, portando il calcolo spaziale da "giocattolo per sviluppatori" a "primi adopter di massa". Allora l'iPhone potrebbe diventare il compagno di calcolo e il centro dati degli occhiali Vision.
  \item \textbf{Occhiali AR}: aprono una "seconda scrivania" davanti ai tuoi occhi, le informazioni digitali diventano onnipresenti come l'aria. Gli occhiali intelligenti dalla partnership tra Meta e Ray-Ban hanno venduto milioni di pezzi, e in Cina aziende come Rokid, XREAL, INMO stanno lanciando occhiali AR di consumo. Gli occhiali AR del 2025-2026 già fanno traduzione in tempo reale, HUD di navigazione e dialogo con assistenti AI, ma sono ancora limitati da display a guida d'onda e autonomia; mancano 3-5 anni per "sostituire il telefono".
  \item \textbf{Anelli intelligenti e indossabili}: Samsung Galaxy Ring, Oura Ring e altri anelli intelligenti stanno diventando popolari, come "sensori corporei" del telefono che integrano il monitoraggio della salute. Il ruolo del telefono sta passando da "terminale universale" a combinazione di "hub + accessori leggeri multipli".
\end{itemize}

\booksection{Interfacce cervello-computer: il pensiero come mezzo di interazione}

Il pensiero stesso diventerà mezzo di interazione, ma questa strada è più lunga del previsto:

\begin{itemize}
  \item \textbf{Ultimi sviluppi di Neuralink}: Neuralink di Musk ha completato nel 2024 il primo impianto su umano; il primo volontario, Noland Arbaugh, controlla col pensiero il mouse del computer per giocare a scacchi. Attualmente le BCI (interfacce cervello-computer) impiantate si rivolgono soprattutto a scenari di assistenza per paralitici, con l'obiettivo di aiutarli a "operare" di nuovo il mondo digitale.
  \item \textbf{BCI non invasive}: diverse aziende sviluppano cerchietti non invasivi (come Kernel, OpenBCI), che leggono segnali cerebrali senza chirurgia, ma con precisione di gran lunga inferiore a quella invasiva. Manca ancora molto per "scrivere col pensiero".
  \item \textbf{Prospettive consumer}: le interfacce cervello-computer sono a 10-20 anni dall'interazione consumer con il telefono, ma a breve termine potrebbero essere usate prima in accessibilità, monitoraggio della concentrazione e altri scenari. Le questioni etiche e di privacy ("i miei segnali cerebrali di chi sono come privacy?") diventeranno motivo di dibattito prima ancora della maturazione tecnologica.
\end{itemize}

\booksection{Il telefono “invisibile” e la nuova pietra angolare dell'identità}

Un'altra direzione dell'evoluzione futura del telefono potrebbe essere “invisibile” — trasformandosi in un “regista dietro le quinte” onnipresente, a diventare l'unica pietra angolare della tua identità digitale:

\begin{itemize}
  \item \textbf{Autenticazione a due fattori in cloud}: completa in un istante la verifica "tu sei tu", per accedere a qualsiasi servizio senza inserire manualmente codici di verifica.
  \item \textbf{Eredità digitale come impostazione predefinita}: all'attivazione di un nuovo telefono è preimpostato un piano di successione dei beni digitali, così che "ciò che viene dopo" diventi default di fabbrica.
  \item \textbf{Da "centrato sul dispositivo" a "centrato sull'identità"}: la vera cassaforte centrale è il tuo archivio personale di dati in cloud, l'hardware è solo una porta di accesso. Cambi telefono, cambi dispositivo, ma identità e dati ti seguono sempre.
\end{itemize}

\booksection{Riflessioni etiche: siamo ancora i padroni del telefono?}

Un interrogativo ultimo resta sempre: siamo “dei” potenziati dalla tecnologia, o un “cervello in una vasca” nutrito da algoritmi?

\begin{itemize}
  \item \textbf{Etica dell'AI}: quando l'AI sul dispositivo capisce le tue preferenze meglio di te, la tua "scelta" è volontà autonoma o un risultato indotto? Quando un AI Agent ordina o risponde ai messaggi per te, dove sono responsabilità e confini?
  \item \textbf{Sovranità dei dati}: i tuoi dati sanitari, i segnali cerebrali, i dati biometrici, a chi appartengono — a te, al produttore del dispositivo o al fornitore cloud? La "Legge sulla protezione delle informazioni personali" cinese e il GDPR dell'UE sono solo l'inizio; l'effettiva attuazione del "diritto alla portabilità dei dati" e del "diritto all'oblio" è ancora lunga.
  \item \textbf{Divario digitale}: telefoni AI, schermi pieghevoli, comunicazione satellitare stanno allargando il divario tra "chi può permetterseli" e "chi non può". L'equo accesso alle infrastrutture digitali diventerà una nuova questione sociale — critica come 20 anni fa "saper navigare o meno".
  \item Quando l'algoritmo sa cosa vuoi vedere meglio di te, quanto del tuo "piacere" è spontaneo?
  \item Quando accendi lo schermo inavvertitamente ogni pochi minuti, è un'abitudine tua o un "comportamento di dipendenza" progettato?
\end{itemize}

\booksection{Restare connessi, ma anche indipendenti — la tua filosofia di vita digitale}

Qualsiasi evoluzione tecnologica, queste filosofie di vita digitale rimarranno valide anche tra 10 anni:

\begin{enumerate}
  \item Sii allocatore di risorse, non ricevitore passivo di notifiche.
  \item Rendi l'automazione il tuo fedele "copilota".
  \item Costruisci il tuo solido "retroterra digitale".
  \item Esegui periodicamente "decluttering digitale" e "digiuno sociale".
  \item Considera sempre sicurezza e privacy come prerequisito di ogni comodità.
\end{enumerate}

\begin{tcbpractice}{Prova: redigi il tuo "Termini di servizio per l'uso di dispositivi digitali personali"}
Scrivi un semplice “termini di servizio” tra te e il tuo telefono (e i futuri occhiali AR, assistenti AI), per definire la vostra relazione. Ad esempio:

\begin{itemize}
  \item Articolo 1: il mio tempo di attenzione mi appartiene. Quando lavoro o sto con la famiglia devi entrare attivamente in modalità "non disturbare".
  \item Articolo 2: i miei dati personali sono una mia proprietà. La tua prima responsabilità è la crittografia sicura e la sincronizzazione comoda.
  \item Articolo 3: il nostro obiettivo comune è migliorare l'efficienza e l'esperienza di vita. Se con design persuasivi mi fai assuefare, ho il diritto di disinstallarti.
\end{itemize}

\end{tcbpractice}

\begin{tcbsummary}{Riepilogo del capitolo}
\begin{itemize}
  \item Vision Pro e occhiali AR sono nuove porte di calcolo oltre il telefono, ma a breve termine difficilmente lo sostituiranno
  \item Le interfacce cervello-computer sono a 10-20 anni dal consumer; le questioni etiche precedono la maturazione tecnologica
  \item Il telefono si evolve verso l'"invisibilità", da centro del dispositivo a centro dell'identità
  \item Etica dell'AI, sovranità dei dati e divario digitale sono i temi centrali dei prossimi 10 anni
  \item Restare connessi ma anche indipendenti è la filosofia immutabile della vita digitale
\end{itemize}

\end{tcbsummary}
